% =====================================================================
% Strand Three: epistemic actions and introspective access.
% Regime 6: the model may query lumped quotients (epimorphic images)
% of the World HMM. The analysis separates genuine introspective
% access to the model's own belief state from first-order competence
% that merely uses it.
% =====================================================================

\subsection*{Epistemic actions: queryable quotients and introspective access}

The regimes above vary \emph{who} acts; this strand varies \emph{what acting is
for}. Following the classical distinction of \citep{kirsh1994epistemic}, the
SWITCH actions of Regimes 2--5 are simultaneously \emph{pragmatic} (they
transition the World state) and \emph{epistemic} (the response is informative
about it). Regime 6 adds a channel that is purely epistemic: it yields
information about $w$ without disturbing it.

\paragraph{Setup (Regime 6).} Fix a family of epimorphisms
$\varphi_i : W \to W^{(i)}$ onto smaller HMMs, chosen so that each quotient is a
lumpable coarse-graining of the World chain \citep{kemeny1960finite}. At
designated positions the model may emit a query token $q_i$; the environment
answers with the current coarse state $y = \varphi_i(w_t)$, leaving $w_t$
untouched. Queries are budgeted (a fixed number per trajectory, or a small
reward cost), so the model must choose \emph{which} quotient to consult. The
cleanest instantiation factors the World outright, $W \cong A \times B$, with
the $\varphi_i$ the two projections: the model holding belief $b_W$ must decide
whether it is more ignorant about the $A$ factor or the $B$ factor. The
intuition motivating the regime is that spending the budget well requires the
model to \emph{know what it knows}: a metacognitive judgement about its own
belief state, made before any new evidence arrives.

\paragraph{The deflationary observation.} We state at the outset the objection
that this intuition does not survive unmodified, because the objection shapes
the entire analysis plan. For an ideal Bayes-adaptive agent the expected
information gain of query $i$ is the entropy of the pushforward of its current
belief, $\mathrm{EIG}_i = H(\varphi_{i\,\sharp}\, b_W)$: the answer is a
deterministic function of $w$, so the mutual information $I(w;\varphi_i(w))$
reduces to the entropy of the coarse marginal. The optimal query policy
$\arg\max_i H(\varphi_{i\,\sharp}\, b_W)$ is therefore a functional of the
first-order belief --- one more readout of $b_W$, requiring no representation
\emph{of} the belief over and above the belief itself. Optimal query behaviour
alone is thus compatible with a pure world model. This is precisely the
deflationary critique long aimed at animal and infant metacognition: opt-out
and help-seeking paradigms \citep{hampton2001monkeys,goupil2016infants} admit
first-order reconstructions in which an uncertainty signal is \emph{used}
without being \emph{represented} \citep{carruthers2008skeptical}, and
behavioural evidence alone has not resolved that dispute. The toy setting
cannot evade the equivalence at the behavioural level --- but it can break it
at the mechanistic level, in two ways.

\paragraph{Escape one: dissociation under miscalibration.} The behavioural
equivalence holds only for the ideal agent, for whom ``my uncertainty'' is a
deterministic function of the observable history. A finite trained model's
realised belief $\hat b_W$ --- the object our probes decode --- deviates from
the analytic posterior $b_W$ somewhere: off the training distribution, on long
or ambiguous histories, at capacity limits. There the hypotheses come apart.
\emph{Introspective access} predicts queries that track the model's own
miscalibration, $\arg\max_i H(\varphi_{i\,\sharp}\, \hat b_W)$: the model asks
about what \emph{it} is uncertain of, even where an ideal observer would not
be. A \emph{compiled heuristic} --- a query policy fit to the typical
uncertainty profile of surface history features --- instead approximates the
ideal policy on-distribution and tracks neither posterior where they diverge.
Because both $b_W$ (analytically) and $\hat b_W$ (by the belief-geometry
probes) are available to us, we can locate histories on which they disagree and
classify every query the model makes. Knowing-what-you-know, on this
operationalisation, is not being right about the world; it is being right about
one's own errors.

\paragraph{Escape two: causal introspection.} The sharper test is
interventional and reuses machinery this proposal already builds. Once the
World-belief subspace is identified, we edit $\hat b_W$ in the residual stream
--- writing in a belief with a different marginal-entropy profile --- and ask
whether the query distribution shifts as the edited belief predicts. A positive
result establishes a causal pathway from the belief representation to the query
choice: the metacognitive act \emph{reads} the model's own internal state
rather than recomputing a shortcut from the history. This is the toy-model
analogue of concept-injection tests of introspection in language models, and
connects directly to the on-policy-detection findings in our forthcoming work
with Lindsey. We note the regress this stops: one can always redescribe
``reading one's own belief state'' as more first-order computation, but at that
point the deflation has no further content --- a circuit that consumes the
belief representation \emph{as} a representation, tracks its idiosyncratic
errors, and is causally load-bearing for behaviour is everything introspective
access could be in any physical system.

\paragraph{The epistemic self and its geometry.} If the introspective readout
exists, Regime 6 adds a third notion of self to the two already in play
(token-type and actor-identity): an \emph{epistemic self}, a represented
quantity whose content is a property of the model's own belief state rather
than of the World. In the factored instantiation the predicted object is
concrete: scalar readouts of the marginal entropies $H(b_A)$ and $H(b_B)$,
linearly available upstream of the query decision. These quantities sit on the
Self side of the Markov blanket while being \emph{about} the World --- the
first genuinely second-order representation the SWITCH demands, and the
self-directed counterpart of the other-directed belief fibers of the previous
strand. This symmetry raises a sharp circuit-level question, parallel to the
privileged-self-pathway question raised there: does the self-uncertainty
readout route through the same machinery that (with stateful actors) represents
\emph{other} agents' beliefs, or through a privileged direct pathway? The
former is the mechanistic form of the thesis that metacognition is self-directed
mindreading \citep{carruthers2009mindreading}; the SWITCH with both stateful
actors and queryable quotients makes that long-standing question empirical.

\paragraph{Design notes.} (i) Query timing is curriculum-staged: fixed query
slots first, so the model chooses only \emph{which} quotient to consult, with
learned timing and explicit query costs as a second phase. This keeps the
regime free of the cold-start problem that makes open-ended reasoning training
expensive. (ii) The quotient family should be designed so that the identity of
the most informative quotient varies across histories (the product projections
already guarantee this, as uncertainty migrates between factors); otherwise a
static preference over oracles is optimal and the regime tests nothing.
(iii) The regime reuses the Regime 2b training infrastructure, and its analytic
comparators --- filtered posterior, pushforward entropies, value of information
--- are cheap to compute. (iv) A natural extension, deferred on cost grounds,
replaces the external oracle with free-form internal rollouts in which the
model generates both halves of a simulated exchange (the chain-of-thought
analogue); the quotient-oracle version isolates the metacognitive question
without first having to train reasoning from scratch.
